% Generated draft: requires mathematical and source review.
\documentclass[11pt]{article}
\usepackage[margin=25mm]{geometry}
\usepackage{fontspec}
\setmainfont{Latin Modern Roman}
\usepackage{amsmath,amssymb,mathtools,unicode-math}
\setmathfont{Latin Modern Math}
\usepackage{enumitem,needspace,xurl,hyperref}
\setlength{\parindent}{0pt}
\setlength{\parskip}{5pt}
\setlength{\emergencystretch}{3em}
\newcommand{\N}{\mathbb{N}}
\newcommand{\Z}{\mathbb{Z}}
\newcommand{\Q}{\mathbb{Q}}
\newcommand{\R}{\mathbb{R}}
\newcommand{\C}{\mathbb{C}}
\newcommand{\F}{\mathbb{F}}
\newcommand{\A}{\mathbb{A}}
\newcommand{\PP}{\mathbb{P}}
\newcommand{\E}{\mathbb{E}}
\newcommand{\eps}{\varepsilon}
\newcommand{\dd}{\,\mathrm{d}}
\begin{document}
\section*{UnsolvedMath ID: 30002111}
\textbf{Source problem code:} catalogue entry
\subsection{Definitions and mathematical statement}
A compact symplectic toric manifold is, in the standard convention, a compact connected symplectic manifold \((M^{2n},\omega)\) without boundary, equipped with an effective Hamiltonian action of an \(n\)-dimensional torus \(T\). A moment map \(\mu:M\to\mathfrak t^*\), where \(\mathfrak t\) is the Lie algebra of \(T\), has image \(P=\mu(M)\), called its moment or Delzant polytope. Choosing lattice coordinates identifies \(\mathfrak t^*\) with \(\R^n\). Changing the normalization of \(\mu\) translates \(P\); changing the integral basis of the torus lattice changes its coordinates by an element of \(\operatorname{GL}(n,\Z)\).

Let \(g\) be the metric described in the source as a ``toric metric.'' Torus invariance means that the action preserves \(g\). The source does not specify whether ``toric metric'' additionally requires a compatible toric K\"ahler structure: an integrable complex structure \(J\) preserved by the torus, with \(g(X,Y)=\omega(X,JY)\). This distinction affects the class of objects quantified over and remains unresolved.

Write \(\operatorname{Spec}(\Delta_g)\) for the ordinary spectrum of the nonnegative scalar Laplace--Beltrami operator on functions, including eigenvalue multiplicities. No boundary condition is needed under the standard convention above. This spectral datum does not include labels describing the torus representations on eigenspaces.

\textbf{Question.} Does \(\operatorname{Spec}(\Delta_g)\) determine the moment polytope of a compact symplectic toric manifold with toric metric, and consequently determine its symplectomorphism type?

A coordinate-independent formulation of the polytope question, subject to the unresolved meaning of ``toric metric,'' is the following: for every dimension \(2n\) and every two such data sets \((M_i,\omega_i,T_i,\mu_i,g_i)\), \(i=1,2\), does
\[
\operatorname{Spec}(\Delta_{g_1})=\operatorname{Spec}(\Delta_{g_2})
\quad\Longrightarrow\quad
\mu_2(M_2)=A\mu_1(M_1)+b
\]
hold for some \(A\in\operatorname{GL}(n,\Z)\) and \(b\in\R^n\), after choosing integral lattice coordinates? If the acting torus and its identification are fixed, the usual moment-map normalization ambiguity is translation alone. The supplied statement does not specify which equivalence relation on polytopes is intended; the displayed implication records the natural formulation allowing changes of torus lattice coordinates, rather than resolving that omission.

The consequent question about symplectomorphism type asks whether isospectral data in the intended class must admit a diffeomorphism \(\phi:M_1\to M_2\) satisfying \(\phi^*\omega_2=\omega_1\). This conclusion concerns the symplectic forms, which are implicit in the source notation \((M,g)\). Whether the intended conclusion also retains the torus action must be clarified. The metric compatibility and polytope equivalence issues require human mathematical review before a single fully specified conjecture can be extracted.
\subsection{Short English statement}
Can the ordinary Laplace spectrum, including multiplicities, recover the moment polytope and symplectomorphism type of a compact symplectic toric manifold with toric metric? The source leaves the precise metric compatibility assumptions and the equivalence allowed between recovered polytopes unspecified.
\subsection{Sources}
\begin{itemize}
\item Supplied local UnsolvedMath registry record, permanent ID \(30002111\), title ``Spectral Determination of Toric Moment Polytopes,'' exact source code \texttt{catalogue entry}, in the Oberwolfach Reports Open Problems collection. The proposed year is recorded as 2012; no individual proposer is supplied.
\item Supplied report citation: \emph{Geometric Aspects of Spectral Theory (2013)}, Oberwolfach Reports, pp.~2013--2076, DOI \texttt{10.4171/owr/2012/33}. Supplied source link: \url{https://doi.org/10.4171/owr/2012/33}. The supplied background describes these reports as records of mathematical workshops and their open-problem discussions.
\item Supplied literature source, ``OWR report'': \url{https://ems.press/content/serial-article-files/46403}.
\item Supplied literature source, ``Dryden Guillemin Sena-Dias'': \url{https://www.ams.org/journals/tran/2012-364-02/S0002-9947-2011-05412-7/}. The supplied assessment attributes results on equivariant and generic toric variants and line-bundle cases to Dryden, Guillemin, and Sena-Dias; precise theorem statements and hypotheses are not supplied.
\item Literature assessment dated exactly \texttt{2026-08-21}: the registry labels the problem ``partially solved'' and states that the ordinary-spectrum question remains open while the variants just described have been proved. This is supplied dated curation derived from the OWR material, not an independently verified assessment. The supplied background explicitly requests independent verification of time-sensitive status and cited claims.
\item The registry labels the statement \texttt{corrected\_verified} and reports that the official report was checked and confirmed the complete statement, with only a mathematical-delimiter correction to the manifold and metric pair. This verification history is supplied metadata; no independent source consultation is asserted here.
\end{itemize}
\end{document}
